%%=============================================================================
%% Databank en persistentie
%%=============================================================================

\chapter{Databank en persistentie}%
\label{ch:databank-persistentie}

Dit hoofdstuk bespreekt hoe data wordt bewaard. Forcebit Ticketing gebruikt een MySQL-databank met Entity Framework Core als Object-Relational Mapper (ORM). EF Core werkt met C\# entiteiten en vertaalt queries naar databankbewerkingen, zodat de applicatie met objecten kan werken in plaats van rechtstreeks met SQL-tabellen~\autocite{MicrosoftEFCore2026}.

\section{Datamodel}

Het datamodel draait rond vier hoofdtabellen: gebruikers, tickets, berichten en bijlagen. Een gebruiker kan meerdere tickets hebben. Een ticket heeft meerdere berichten. Bijlagen horen bij een ticket en kunnen ook aan een specifiek bericht gekoppeld worden.

\begin{figure}[H]
  \centering
  \vspace{-0.35\baselineskip}
  \includegraphics[width=0.95\linewidth,height=0.33\textheight,keepaspectratio]{Database_schema.png}
  \vspace{-0.35\baselineskip}
  \caption[Databankschema]{Databankschema van de belangrijkste tabellen en relaties.}
  \label{fig:database-schema}
\end{figure}

\clearpage

\section{DbContext en configuraties}

\texttt{AppDbContext} is de centrale EF Core klasse. De \texttt{DbSet} properties tonen welke tabellen de applicatie gebruikt. De detailconfiguratie staat in aparte configuratieklassen, zodat de context kort blijft en elke tabelmapping apart onderhoudbaar is.

\begin{listing}[htbp]
\begin{minted}[fontsize=\scriptsize,baselinestretch=0.9]{csharp}
public DbSet<User> Users => Set<User>();
public DbSet<Ticket> Tickets => Set<Ticket>();
public DbSet<TicketMessage> TicketMessages => Set<TicketMessage>();
public DbSet<TicketAttachment> TicketAttachments => Set<TicketAttachment>();

protected override void OnModelCreating(ModelBuilder modelBuilder)
{
    base.OnModelCreating(modelBuilder);
    modelBuilder.ApplyConfigurationsFromAssembly(
        typeof(AppDbContext).Assembly);
}
\end{minted}
\caption[DbContext met configuraties]{\texttt{AppDbContext} laadt configuraties uit de persistence-laag.}
\label{lst:appdbcontext}
\end{listing}

\section{Enumwaarden en relaties}

Ticketstatus, categorie en onderwerp zijn enums in de domain-laag. Een enum groepeert vaste, benoemde waarden binnen een eigen type~\autocite{MicrosoftCSharpEnums2026}. In de databank worden deze waarden als tekst opgeslagen. Dat maakt de inhoud leesbaarder tijdens testen en controle. Een rij met status \texttt{Open} is duidelijker dan een numerieke waarde zoals \texttt{1}.

De relatie tussen ticket en klant gebruikt \texttt{DeleteBehavior.Restrict}. Daardoor kan een gebruiker niet zomaar verwijderd worden als tickets nog naar die gebruiker verwijzen. Voor een supportapplicatie is dat logisch, omdat ticketgeschiedenis belangrijk blijft.

\section{Repository pattern}

Repositories schermen querydetails af van de service-laag. \texttt{TicketService} vraagt bijvoorbeeld een ticketdetail op, maar hoeft niet te weten welke relaties mee geladen worden. Dat maakt de service beter leesbaar: de service beschrijft de workflow, de repository beschrijft de databankquery.
